\section*{Notation index}
\addcontentsline{toc}{section}{Notation index}
\label{intro:notation}

This index collects objects defined in the text. The one-dimensional
model uses $x\in\T_{2\pi}$, a fixed parameter $0<d<1/2$, and
$\omega_d(x)=\sqrt{1-2d\cos x}$. In that section, $\Xi$ is the coarea
measure on the regular resonance surface and
$m(\dd x)=\dd x/(2\pi)$ is normalized circle measure.
The table below mainly concerns the three-dimensional model.

\begingroup
\small
\renewcommand{\arraystretch}{1.25}
\begin{longtable}{@{}p{0.27\linewidth}p{0.69\linewidth}@{}}
\toprule
Notation & Meaning and convention\\
\midrule
\endhead
$k\in D=[-R,R]^3$, $X\in\T^3$
 & $R>0$ is fixed; $k^{(j)}$ is a momentum component and $k_i$ a collision
   leg. The velocity is $v(k)=2k$ and $\T^3=\T_{2\pi}^3$.
   The measures $\dd X$ and $\dd k$ are periodic and ordinary
   Lebesgue measure, respectively.\\
$\Delta\varphi$, $\Xi_D$, $\Gamma_k$
 & $\Delta\varphi=\varphi_0+\varphi_1-\varphi_2-\varphi_3$,
   $\varphi_i=\varphi(k_i)$. The four-leg resonance measure is $\Xi_D$;
   $\Gamma_k$ is its fibre measure with output leg $k_0=k$ fixed.\\
$\mathcal T$, $c=4\pi/\varepsilon$
 & $\mathcal T=\partial_t+v\cdot\nabla_X$.
   The kinetic equation is $\mathcal Tf_c=c\mathcal C(f_c)$, and
   the limit is $c\to\infty$.\\
$\Psi$, $\Theta_+=\mathcal O$
 & $\Psi=(1,k^{(1)},k^{(2)},k^{(3)},|k|^2)^{\mathsf T}$;
   $\Theta_+=\{\theta:\min_D\theta\cdot\Psi>0\}$ is the positive
   parameter domain.\\
$N_\theta$, $\bar N$, $N_c$
 & $N_\theta=(\theta\cdot\Psi)^{-1}$.
   The prescribed Euler equilibrium is $\bar N=N_{\bar\theta}$;
   $N_c=N_{\theta_c}$ matches all five moments of $f_c$.
   Each local abbreviation $N$ is specified where it is used.\\
$\kappa_N=N\Psi$, $P_N,Q_N$
 & $P_N$ is the unweighted $L^2$ orthogonal projection onto
   $\Span\{N\Psi_a\}_{a=0}^4$, extended to the weighted space;
   $Q_N=I-P_N$.\\
$\mathscr H(f\mid N)$, $H_c(t)$
 & Relative entropy is defined in \eqref{intro:relative-entropy}.
   Here $H_c=\int_{\T^3}\mathscr H(f_c\mid N_c)\dd X$.
   Pointwise moment matching gives the exact decomposition
   \eqref{hyd:entropy-pythagoras}.\\
$\nu_*$, $H_\nu$
 & $\nu_*$ is the collision frequency of a fixed positive RJ state,
   and $H_\nu=L^2(D,\nu_*\dd k)$.
   On compact positive parameter sets, $\nu_N\asymp\nu_*$.\\
$Z$, $r(k)$, $E_\eta(\xi)$
 & $Z=\{\pm R\}^3$ is the corner set and $r(k)=\dist_\infty(k,Z)$.
   Set $E_\eta(\xi)=\{k\in D:\|k-\xi\|_\infty\le\eta\}$ and
   $E_\eta=\bigcup_{\xi\in Z}E_\eta(\xi)$.\\
$\mathcal X$, $\mathcal Y$
 & $\mathcal X=L^\infty(D)$ and
   $\mathcal Y=\{h:\nu_*h\in L^\infty(D)\}$,
   with norms $\|h\|_\infty$ and $\|\nu_*h\|_\infty$.\\
$L_N$, $\mathfrak a_N$
 & $L_N=-N^{-1}D\mathcal C(N)N$.
   The collision form $\mathfrak a_N$ is \eqref{geom:form};
   $L_{N,Q}^{-1}$ is its weighted inverse with five moment
   normalizations.\\
$q_c,G_c$, $e_c,a$
 & $q_c$ solves the nonlinear resolvent equation and
   $G_c=\bar Nq_c$.
   The errors are $e_c=f_c-N_c-G_c$ and $a=\theta_c-\bar\theta$.
   Section \ref{hyd:section} abbreviates $e_c,G_c$ as $e,G$
   at fixed $c$. Its adapted derivative variables are
   $\eta_\alpha=\partial_X^\alpha e/N_c$.\\
$\xi=(g,B,h)$, $B^\circ$
 & For $\theta=(\alpha,\beta,\gamma)$, set $g=\nabla_X\alpha$,
   $B_{jl}=\partial_{X_j}\beta_l$, $h=\nabla_X\gamma$, and
   $B^\circ=(B+B^{\mathsf T})/2-(\tr B)I/3$.\\
$\mathsf K_{ij}^{ab}$, $\mathsf D_N(\ell)$
 & The Onsager tensor and spatial Fourier dissipation symbol,
   defined by \eqref{ons:definition} and \eqref{ons:diffusion-symbol}.
   The first-order coefficient is $c^{-1}\mathsf K$.\\
$z_c$, $u$, $\mathsf R_N$
 & $z_c=c(1-N_c/f_c)$.
   Its limit satisfies $L_{\bar N}u=-\mathcal T\log\bar N$,
   $P_{\bar N}u=0$.
   The operator $\mathsf R_Nh=\tfrac12\Delta(h/N)$ takes values
   in the common resonance space.\\
$\mathcal J_c$, $Z_{c,i}^a$
 & $\mathcal J_c=\tfrac c2\sqrt{\prod_i f_{c,i}}\Delta(1/f_c)$
   is the signed resonance current.
   The residual flux of quantity $a$ in direction $i$ is
   $Z_{c,i}^a=\int_Dv_i\Psi_a(f_c-N_c)\dd k$.\\
\bottomrule
\end{longtable}
\endgroup

For an integer $s\ge0$, the squared $H_X^sH_\nu$ norm is
$\sum_{|\alpha|\le s}\int_{\T^3\times D}
\nu_*|\partial_X^\alpha h|^2\dd X\dd k$.
The space $H^{-1}(\T^3)$ is the dual of $H^1(\T^3)$.
Time norms use the stated interval; time-integrated convergence
does not include strong endpoint traces.
Momentum integrals without an explicit domain use the domain of
the model under discussion.
